\paragraph{\trConsistent{}地选取\trPerturbationDatum。}\label{sub:consistent-perturbation-data}
这里我们需要简要地重访\trGluingConstruction。设有一棵 $d$ \trLeafedTree{}$T$，并且对每个\trVertex{}$v\in\mathrm{Ve}(T)$，有一个带\trStriplikeEnd{}与\trLagrange\trLabel{}的 $|v|$ \trPointedDisc{}$S_v$。回顾：在诸\trFlag{}$f=(v,e)$ 与 $S_v$ 的诸无穷远点之间存在一个\trBijective。于是，在我们的处境中，每个\trFlag{}都关联着一对\trExact\trLagrange\trSubmanifold。我们称诸 $S_v$ 是\textbf{\trLabel\trCompatible}（compatibly labeled）的，若对每条 $e\in\mathrm{Ed}^{\mathrm{int}}(T)$，两个\trFlag{}$f^{\pm}(e)=(v^{\pm}(e),e)$ 所关联的\trLagrange\trSubmanifold{}对相同。要直观地想象这个局面，可以用\trExact\trLagrange\trSubmanifold{}给 $\RR^2\setminus T$ 的诸\trConnectedComponent{}贴上\trLabel（见\cref{fig:lagrangian-label-compatibility}）。把诸 $S_v$ \trGluing{}起来所得的 $(d{+}1)$ \trPointedDisc\trFamily{}$\scr{S}$，便以显然的方式承袭了\trLagrange\trLabel。再设诸 $S_v$ 还配备了\trPerturbationDatum{}$(K_v,J_v)$，使得在每一对端 $\epsilon_{f^{\pm}(e)}$ 上相关联的\trFloerDatum{}都相同。它们便经由\trGluingConstruction{}确定 $\scr{S}$ 上的一个\trPerturbationDatum{}$(K,J)$。这个数据可\trSmooth{}地延拓到自然的\trPartialCompactification{}$\bar{\scr{S}}$ 上，其含义是：存在
\begin{equation}\label{eq:compactified-perturbation-datum}
\bar{K}\in\Omega^1_{\bar{\scr{S}}/\bar{\scr{R}}}(\bar{\scr{S}},\mathscr{H})
\qquad\text{与}\qquad
\bar{J}\in C^\infty(\bar{\scr{S}},\mathscr{J})
\end{equation}
其限制到 $\scr{S}$ 上等于 $(K,J)$（逐\trFiber{}的 1-形式空间在\trPartialCompactification{}中是有意义的，因为 $\bar{\scr{S}}\to\bar{\scr{R}}$ 是一个\trSubmersion）。在最\trDegenerate{}的点 $0\in\bar{\scr{R}}$ 之上——那里\trFiber{}为 $\scr{S}_0=\coprod_v S_v$——$(\bar{K},\bar{J})$ 恰是所给数据 $(K_v,J_v)$ 的\trDisjointUnion。最后，如果我们\trGluing{}的是\trFamily{}$\scr{S}_v$ 而非单个\trSurface，那么以上全部论述依然成立。
取定\trStriplikeEnd{}的一个\trUniversal{}选取。此外，对所考虑的\trSymplectic\trManifold{}$M$ 中每一对\trExact\trLagrange\trSubmanifold{}$(L_0,L_1)$，取定一个\trFloerDatum{}$(H_{L_0L_1},J_{L_0L_1})$。这时，所谓\textbf{\trPerturbationDatum{}的\trUniversal{}选取}（universal choice of perturbation data），是指对每个 $d\ge 2$ 与每个由\trExact\trLagrange\trSubmanifold{}构成的 $(d{+}1)$ 元组 $(L_0,\ldots,L_d)$，选取一个\trPerturbationDatum
\begin{equation}\label{eq:universal-perturbation-data}
(K_{L_0,\ldots,L_d},J_{L_0,\ldots,L_d})
\end{equation}
其定义在带\trLagrange\trLabel{}$(L_0,\ldots,L_d)$ 的\trFamily{}$\scr{S}^{d+1}\to\scr{R}^{d+1}$ 上，并且在该\trFamily{}的诸\trStriplikeEnd{}上\trDegenerate{}为与诸对\eqref{eq:lagrangian-pairs} 相配的\trFloerDatum。
\begin{figure}[htbp]
\centering
\begin{tikzpicture}[line cap=round,line join=round,scale=0.88]
\begin{scope}[shift={(0,0)}]
  \draw (-1.7,0)--(0,0);
  \draw (0,0)--(1.05,0.8);
  \draw (1.05,0.8)--(2.15,1.4);
  \draw (1.05,0.8)--(2.25,0.5);
  \draw (0,0)--(1.95,0);
  \draw (0,0)--(1.6,-0.9);
  \node at (-0.14,-0.4) {$v_1$};
  \node at (0.78,1.06) {$v_2$};
  \node at (0.15,-1.45) {\small 树 $T$};
\end{scope}
\begin{scope}[shift={(5.05,0.25)}]
  \draw[fill=gray!45] (0,0) circle (0.95);
  \fill (180:0.95) circle (2.1pt);
  \fill (58:0.95) circle (2.1pt);
  \fill (-12:0.95) circle (2.1pt);
  \fill (-64:0.95) circle (2.1pt);
  \node at (-1.78,0.32) {$L_{v_1,3}$};
  \node at (1.6,0.62) {$L_{v_1,2}$};
  \node at (1.62,-0.28) {$L_{v_1,1}$};
  \node at (-0.72,-1.32) {$L_{v_1,0}$};
  \begin{scope}[shift={(3.5,0)}]
    \draw[fill=gray!45] (0,0) circle (0.95);
    \fill (96:0.95) circle (2.1pt);
    \fill (-4:0.95) circle (2.1pt);
    \fill (-72:0.95) circle (2.1pt);
    \node at (0,1.42) {$L_{v_2,2}$};
    \node at (1.6,0.02) {$L_{v_2,1}$};
    \node at (0,-1.4) {$L_{v_2,0}$};
  \end{scope}
  \node at (1.75,-2.05) {\small 与 $\mathrm{Ve}(T)$ 相关联的\trPointedDisc，};
  \node at (1.75,-2.5) {\small 及其\trLagrange\trLabel};
\end{scope}
\begin{scope}[shift={(0,-5.4)}]
  \draw (-1.7,0)--(0,0);
  \draw (0,0)--(1.05,0.8);
  \draw (1.05,0.8)--(2.3,1.5);
  \draw (1.05,0.8)--(2.45,0.3);
  \draw (0,0)--(2.0,0);
  \draw (0,0)--(1.75,-1.35);
  \node at (-0.45,1.12) {$L_{v_1,3}=L_{v_2,2}$};
  \node at (2.72,0.88) {$L_{v_2,1}$};
  \node at (2.15,0.16) {$L_{v_1,2}=L_{v_2,0}$};
  \node at (1.45,-0.5) {$L_{v_1,1}$};
  \node at (-0.35,-1.05) {$L_{v_1,0}$};
  \node at (0.35,-1.85) {\small \trCompatible{}性条件};
\end{scope}
\begin{scope}[shift={(5.75,-5.2)}]
  \path[fill=gray!45]
    (-1.0,0.15)
    .. controls (-1.12,0.8) and (-0.5,1.2) .. (0.1,1.1)
    .. controls (0.55,1.02) and (0.62,0.72) .. (1.0,0.66)
    .. controls (1.28,0.62) and (1.35,0.95) .. (1.72,1.05)
    .. controls (2.25,1.18) and (2.8,0.98) .. (2.85,0.5)
    .. controls (2.9,0.1) and (2.55,-0.2) .. (2.2,-0.1)
    .. controls (1.98,0.06) and (1.9,0.2) .. (1.7,0.12)
    .. controls (1.5,0.04) and (1.48,-0.28) .. (1.3,-0.42)
    .. controls (1.0,-0.66) and (0.3,-0.8) .. (-0.2,-0.62)
    .. controls (-0.68,-0.46) and (-0.92,-0.28) .. cycle;
  \fill (-1.0,0.15) circle (2.1pt);
  \fill (2.42,0.94) circle (2.1pt);
  \fill (2.62,-0.02) circle (2.1pt);
  \fill (1.42,-0.2) circle (2.1pt);
  \fill (0.2,-0.71) circle (2.1pt);
  \node[anchor=south east] at (-1.02,0.42) {$L_{v_1,3}$};
  \node at (3.42,0.82) {$L_{v_2,1}$};
  \node at (3.3,-0.42) {$L_{v_2,0}$};
  \node at (2.2,-0.52) {$L_{v_1,1}$};
  \node at (0.16,-1.14) {$L_{v_1,0}$};
  \node at (0.9,-1.75) {\small \trGluing{}成的\trSurface{}及其\trLabel};
\end{scope}
\end{tikzpicture}
\caption{}\label{fig:lagrangian-label-compatibility}
\end{figure}
从现在起，假设我们所取的\trStriplikeEnd{}的\trUniversal{}选取是\trConsistent{}的。设我们已经作了\trPerturbationDatum{}的一个\trUniversal{}选取。令 $T$ 是一棵稳定的 $d$ \trLeafedTree；把诸\trUniversal\trFamily{}$\scr{S}^{|v|}$ 与它的诸\trVertex{}相关联；并为它们配备\trCompatible{}的\trLagrange\trLabel。这时，\trGluing{}而成的\trFamily{}$\scr{S}\to\scr{R}$ 带有两组\trPerturbationDatum，它们\trApriori{}地看可能不同（更确切地说，它们在诸\trStriplikeEnd{}上总相同，但在\trSurface{}的其余部分上可能不同）。第一组是由诸 $\scr{S}^{|v|}$ 上的\trPerturbationDatum{}经\trGluingConstruction{}诱导而来；另一组则是用\trClassifyingMap{}$\Gamma^T$，从 $\scr{S}^{d+1}$ 上（带正确\trLabel）的\trPerturbationDatum{}的\trUniversal{}选取\trPullback{}而来。我们称\trPerturbationDatum{}的这个\trUniversal{}选取是\emph{\trConsistent}的，若在一切这类处境中，下列条件都得到满足：
\begin{itemize}
\item 存在一个子集 $U\subset\scr{R}$，使得\trGluingParameter{}在其中充分小，并且上述两组\trPerturbationDatum{}在诸\trSurface{}$\scr{S}_r$（$r\in U$）的薄部上相同。这里指的是由\trGluing{}过程产生的\trThickThin{}分解，见\cref{thm:thick-thin-decomposition}。
\item 令 $(K,J)$ 为 $\scr{S}$ 上的第一组\trPerturbationDatum{}（由\trGluing{}所得的那组），而 $(\bar{K},\bar{J})$ 为其如\eqref{eq:compactified-perturbation-datum} 那样到\trPartialCompactification{}$\bar{\scr{S}}$ 上的延拓。我们要求另一组数据（由从 $\scr{S}^{d+1}$ \trPullback{}而得者）也能\trSmooth{}地延拓到 $\bar{\scr{S}}$，并且其延拓在子集 $\{0\}^{\mathrm{Ed}^{\mathrm{int}}(T)}\times\scr{R}^T\subset\bar{\scr{R}}$——即一切\trGluingParameter{}皆为零之处——上与 $(\bar{K},\bar{J})$ 相同。
\end{itemize}
为了更好地体会\trConsistency{}的含义，设有一列 $r_1,r_2,\ldots\in\scr{R}^{d+1}$，\trConvergence{}到 $\bar{\scr{R}}^{d+1}\setminus\scr{R}^{d+1}$ 的 $T$ \trStratum{}中的一点。那么对大的 $k$，可以写成
\begin{equation}\label{eq:sequence-near-stratum}
r_k=\gamma^T(\{\rho_{e,k}\},\{r_{v,k}\})
\end{equation}
其中 $r_{v,k}\in\scr{R}^{|v|}$（$v\in\mathrm{Ve}(T)$），而 $\rho_{e,k}\in(-1;0)$（$e\in\mathrm{Ed}^{\mathrm{int}}(T)$）。当 $k\to\infty$ 时，每一列 $r_{v,k}$ 都\trConvergence，而诸 $\rho_{e,k}$ 趋于 $0$。更几何地，考虑诸\trPointedDisc
\begin{equation*}
S_k=\scr{S}^{d+1}_{r_k},\qquad S_{v,k}=\scr{S}^{|v|}_{r_{v,k}}.
\end{equation*}
于是\eqref{eq:sequence-near-stratum} 说的是：若按 $T$ 诸\trEdge{}所指示的那样把诸 $S_{v,k}$ \trGluing{}起来，并在任意两者之间留下有限\trStrip{}$[0;-\log(-\rho_{e,k})/\pi]\times[0;1]$（当 $k\to\infty$ 时越来越长），则所得结果\trIsomorphic{}于 $S_k$。为诸 $S_k$ 配备某些\trLagrange\trLabel{}$(L_0,\ldots,L_d)$，并为诸 $S_{v,k}$ 配备相应的\trConsistent\trLabel。按定义，\trPerturbationDatum{}的一个\trUniversal{}选取在每个带\trLagrange\trLabel{}的稳定\trPointedDisc{}上都给出一个取定的\trPerturbationDatum，于是在我们的情形，便在诸 $S_k$ 与 $S_{v,k}$ 上给出了取定的\trPerturbationDatum。\trConsistency{}说的是：对大的 $k$，$S_k$ 上的\trPerturbationDatum{}“\trConvergence{}到”$\coprod_v S_{v,k}$ 上的那些；这个说法之所以有意义，是因为按定义，$S_k$ 乃是 $\coprod_v S_{v,k}$ 的一个子集的\trQuotient。此外，在 $S_k$ 的诸有限\trStrip{}上（其诸\trStriplikeEnd{}上亦然），这种\trConvergence{}是平凡的：当 $k\gg 0$ 时，那里的\trPerturbationDatum{}将与 $\coprod_v S_{v,k}$ 相应部分上的那些相等。
\begin{lemma}\label{thm:consistent-perturbation-data-exist}
\trPerturbationDatum{}的\trConsistent\trUniversal{}选取是存在的。\noqed
\end{lemma}
其证明与\cref{thm:consistent-ends-exist}属于同一类型的归纳论证；关于这个证明，我们只限于讲一点观察。一般地，一个\trSurface{}$S=\scr{S}^{d_0+1}_r$ 可能落在不止一个\trGluingMap{}$\gamma^T$ 的像之中，而这些\trGluingMap{}给出不同的\trThickThin{}分解。不过，可以对每个 $T$ 选取一个开子集 $U^T\subset(-1;0)^{\mathrm{Ed}^{\mathrm{int}}(T)}\times\scr{R}^T$，使\trGluingParameter{}在其中很小，并且使得对每个 $r\in\scr{R}^{d+1}$，都存在一棵\trMaximalTree{}$T$ 满足 $r\in\gamma^T(U^T)$（极大性的含义是：具有该性质的其余\trTree{}$\bar{T}$ 都可由 $T$ \trCollapse{}诸\trEdge{}而得）。作为\trStriplikeEnd{}选取之\trConsistency{}的一个推论，由 $\gamma^T$ 诱导的 $S$ 的\trThickThin{}分解是极大的（来自其余 $\gamma^{\bar{T}}$——其中 $r\in\gamma^{\bar{T}}(U^{\bar{T}})$——的诸分解，都可由它遗忘掉一些有限薄带而得）。相对于我们的\trConsistency{}要求而言，极大的分解对\trPerturbationDatum{}施加的条件最为严苛。这对于证明\cref{thm:consistent-perturbation-data-exist}所用的策略有如下含义。设我们已经为 $d<d_0$ 选好了\trPerturbationDatum{}\eqref{eq:universal-perturbation-data}，使得\trConsistency{}条件得到满足。那么，对 $\scr{S}^{d_0+1}$ 的构造就从那些最复杂（\trEdge{}数最多）的\trTree{}$T$ 的子集 $\gamma^T(U^T)$ 出发——这些 $T$ 对应于 $\bar{\scr{R}}^{d_0+1}$ 的诸\trVertex——再按\trCompactification{}诸\trBoundaryStratum{}维数的递增次序进行下去。
\begin{remark}\label{thm:stronger-consistency-remark}
初看之下（特别是把上面的讨论与\cref{sub:consistent-striplike-ends}中的相应讨论相比较时），施加一个更强的\trConsistency{}条件似乎更自然，即要求\trGluing{}而成的\trFamily{}$\scr{S}\to\scr{R}$ 上的两组\trPerturbationDatum{}在\trGluingParameter{}取小值时相同。为了看清这种做法的缺陷何在，取一棵带两个\trVertex{}的\trTree{}$T$，其\trValency{}各为 $3$。给每个\trVertex{}关联上\trFamily{}$\scr{S}^3\to\scr{R}^3$，它实际上只由唯一那个 $3$ \trPointedDisc{}$S$ 构成。用\eqref{eq:striplike-gluing} 的记号，所得的\trGluing{} $4$ \trPointedDisc{}是 $S\#_l S$。把所有\trLabel{}都取成同一个\trExact\trLagrange\trSubmanifold。若一个\trUniversal\trPerturbationDatum{}满足更强的\trConsistency{}条件，它便为诸\trDisk{}$S\#_l S$ 配上了在“两半”上都相同的\trPerturbationDatum（更确切地说，在厚部的两个\trConnectedComponent{}上相同）。当试图在这类\trPerturbationDatum{}之内为\trInhomogeneous\trPseudoholomorphicMap{}的诸\trModuliSpace{}达到\trRegularity{}时，就会出现\trMultipleCover{}型的问题。我们对\trConsistency{}的定义绕开了这一困难。
\end{remark}
